% abstract.tex -- Paper VIII
On the A=1 Discrete Causal Lattice the magnetic field enters the tick rule as a
spatial link phase and yields, through the bipartite-plaquette Wilson loop, an exact
induced-action tensor with eigenvalues $\{4,4,16\}$ and optical axis $(1,1,-1)$
(Paper~I, Appendix~B). The electric field enters differently---as an on-site,
mass-like phase advance---so it admits no Wilson-loop analogue, and the corresponding
electric action block (the permittivity $\varepsilon$) is absent from both Paper~I and
Paper~II. Because the framework establishes only the spatial point group $O_h$ and not
the Lorentz boosts, there is no symmetry route from the magnetic $\mu^{-1}$ to the
electric $\varepsilon$: the electric block must be derived directly. We carry out that
derivation and assemble the covariant $(\varepsilon,\mu^{-1})$ completion of the
lattice gauge response. The result decides a binary physical question left open by
Paper~IV: whether the lattice's vacuum birefringence about $(1,1,-1)$ cancels---so the
framework predicts the \emph{absence} of gauge-sector birefringence---or persists as a
genuine, testable signal. The magnetic $\{4,4,16\}$ result is retained throughout as an
exact consistency anchor.
